% !TeX root = ../../Book1.tex
% 纲要标题：### 6.1 直积与直和
% 纲要位置：计划纲要.md，第 277 行。

\section{直积与直和}
\label{b1:sec:0601}

本节的直积、直和，下一节的半直积，以及\cref{b1:thm:split-extension-semidirect}的分裂判别，是本章首先需要掌握的内容。前面几章主要从一个群中寻找子群和商群，现在反过来，用已知的群构造新群。最直接的办法是在各坐标独立运算，这就是直积。因子有无限多个时，是否允许无限多个非单位坐标便成为实质区别；交换群的直和只允许有限多个。本节同时考察这些构造上的同态，以说明它们各自解决什么问题。

% 本节正文按下列原纲要要点撰写。

% 纲要内容（仅作写作计划，尚非教材正文）：
%
% * 有限直积、映射性质与内直积分解
% * 任意群族的直积与限制直积式直和
% * 交换群范畴中的直和
% * 非交换群的有限支集直积是直积型构造，不是群范畴中的余积；后者由第8.3节的自由积承担
%

\subsection{有限直积、映射性质与内直积分解}
\label{b1:subsec:0601:01}

\begin{definition}\label{b1:def:direct-product}
设 \(r\geq1\) 为整数，\(G_1,\ldots,G_r\) 为群。在笛卡尔积 \(G_1\times\cdots\times G_r\) 上逐坐标定义乘法：
\[
(g_1,\ldots,g_r)(h_1,\ldots,h_r)
=(g_1h_1,\ldots,g_rh_r).
\]
将所得群称为 \(G_1,\ldots,G_r\) 的\term[zhiji]{直积}[direct product]。
\end{definition}

结合律逐坐标成立，单位为 \((1,\ldots,1)\)，逆元为 \((g_1^{-1},\ldots,g_r^{-1})\)，所以这确为群。投影 \(\pi_i:\prod_jG_j\rightarrow G_i\) 和把其他坐标设为单位的嵌入 \(\iota_i:G_i\rightarrow\prod_jG_j\) 都是同态。不同坐标嵌入的像逐元素交换，即使每个 \(G_i\) 自身并不交换。

\begin{proposition}[直积的泛性质]\label{b1:prop:product-universal}
设 \(r\geq1\) 为整数，\(G_1,\ldots,G_r\) 为群，\(\pi_i:G_1\times\cdots\times G_r\rightarrow G_i\) 为第 \(i\) 个坐标投影。给定群 \(H\) 以及群同态 \(f_i:H\rightarrow G_i\)（\(1\leq i\leq r\)），存在唯一群同态 \(f:H\rightarrow G_1\times\cdots\times G_r\) 满足 \(\pi_i f=f_i\) 对每个 \(i\) 成立。
\end{proposition}
\begin{proof}
必有 \(f(h)=\bigl(f_1(h),\ldots,f_r(h)\bigr)\)，这保证唯一性，也给出候选映射。由于每个 \(f_i\) 都保持乘法，候选映射也逐坐标保持乘法，因而确实存在。
\end{proof}

这个泛性质刻画的是映入直积的同态：任意一族坐标同态 \(f_i:H\rightarrow G_i\) 都可以组合，并不要求它们之间另有相容条件。方向倒过来，从各因子到 \(K\) 的同态要合并成一个映出有限直积的同态，就必须保留直积内部已有的关系
\[
\iota_i(g)\iota_j(h)=\iota_j(h)\iota_i(g)
\quad(i\ne j,\ g\in G_i,\ h\in G_j).
\]
因此，不同坐标的像必须逐元素交换；这个条件来自原群的乘法关系，而非各因子自身是否交换。

\begin{proposition}[从有限直积出发的同态]\label{b1:prop:finite-product-maps}
设 \(r\geq1\) 为整数，\(G_1,\ldots,G_r,K\) 为群。令 \(\iota_i:G_i\rightarrow G_1\times\cdots\times G_r\) 为第 \(i\) 个坐标嵌入，并给定群同态 \(u_i:G_i\rightarrow K\)。存在群同态
\[
u:G_1\times\cdots\times G_r\longrightarrow K
\]
满足 \(u\iota_i=u_i\)，当且仅当不同 \(u_i\) 的像逐元素交换，即
\[
u_i(g)u_j(h)=u_j(h)u_i(g)
\quad(i\neq j,\ g\in G_i,\ h\in G_j).
\]
存在时，该同态唯一，且由
\[
u(g_1,\ldots,g_r)=u_1(g_1)\cdots u_r(g_r)
\]
给出。
\end{proposition}
\begin{proof}
若这样的 \(u\) 存在，因不同坐标嵌入的像逐元素交换，对
\(\iota_i(g)\iota_j(h)=\iota_j(h)\iota_i(g)\) 应用 \(u\)，即得到所述交换条件。

反过来，设这些交换条件成立，按给定公式定义 \(u\)。任取 \(g=(g_1,\ldots,g_r)\)、\(h=(h_1,\ldots,h_r)\)，由各 \(u_i\) 的同态性，有
\[
\begin{aligned}
u(gh)
&=u_1(g_1)u_1(h_1)\cdots u_r(g_r)u_r(h_r)\\
&=u_1(g_1)\cdots u_r(g_r)\,u_1(h_1)\cdots u_r(h_r)\\
&=u(g)u(h).
\end{aligned}
\]
中间一步只交换不同坐标的像，故 \(u\) 是同态。又因每个 \(u_j(1)=1\)，把 \(g_i\) 嵌入第 \(i\) 个坐标后，上式的定义给出 \(u\iota_i(g_i)=u_i(g_i)\)。

最后，每个元素均可写成
\((g_1,\ldots,g_r)=\iota_1(g_1)\cdots\iota_r(g_r)\)。任何满足限制条件的同态都必须将它送到 \(u_1(g_1)\cdots u_r(g_r)\)，所以所求同态唯一。
\end{proof}

\begin{proposition}[内直积判别]\label{b1:prop:internal-direct-product}
设 \(G\) 为群，\(N,H\trianglelefteq G\)。若 \(N\cap H=\{1_G\}\)、\(NH=G\)，则称 \(G=NH\) 为 \(G\) 关于 \(N,H\) 的一个\term[nei-zhiji]{内直积分解}[internal direct product decomposition]；此时
\[
N\times H\longrightarrow G,\qquad(n,h)\longmapsto nh
\]
是同构。
\end{proposition}
\begin{proof}
先证明两个因子逐元素交换。对 \(n\in N,h\in H\)，考察它们的交换子
\[
[n,h]\coloneqq nhn^{-1}h^{-1}.
\]
两个子群的正规性分别给出
\[
nhn^{-1}h^{-1}=n(hn^{-1}h^{-1})\in N,
\qquad
nhn^{-1}h^{-1}=(nhn^{-1})h^{-1}\in H.
\]
因此 \([n,h]\in N\cap H=\{1\}\)，即 \([n,h]=1\)，也就是 \(nh=hn\)。

记所给映射为 \(F\)，则
\[
F(n,h)F(n',h')=nhn'h'=nn'hh'=F(nn',hh'),
\]
所以它是同态；\(NH=G\) 给出满射。若 \(F(n,h)=nh=1\)，则 \(n=h^{-1}\in N\cap H\)，故 \(n=h=1\)，核平凡。\cref{b1:prop:kernel-fibers}说明 \(F\) 单射，于是 \(F\) 为同构。
\end{proof}

这些条件分别控制表达式和乘法。\(NH=G\) 保证每个群元素都能写成 \(nh\)；\(N\cap H=\{1\}\) 保证这种表达唯一，因为
\[
nh=n'h'\quad\Longrightarrow\quad n'^{-1}n=h'h^{-1}\in N\cap H,
\]
所以 \(n=n'\)、\(h=h'\)。两个因子都正规，则使交换子同时落入两个因子；再结合交平凡，便得到逐元素交换，使这些坐标的乘法成为直积乘法。只要求一个因子正规时，通常需要下一节的半直积。

\subsection{任意直积与有限支集直积}
\label{b1:subsec:0601:02}

对任意指标集合 \(I\) 和群族 \((G_i)_{i\in I}\)，直积
\(\prod_{i\in I}G_i\) 是全部族 \((g_i)_{i\in I}\) 的集合，仍逐坐标运算。结合律仍在每个坐标独立成立，单位族为 \((1)_{i\in I}\)，逆元为 \((g_i^{-1})_{i\in I}\)，所以这也是群。\cref{b1:prop:product-universal}所述进入直积的泛性质也推广到这里：给定群 \(H\) 与每个 \(i\in I\) 上的同态 \(f_i:H\rightarrow G_i\)，令 \(f(h)=\bigl(f_i(h)\bigr)_{i\in I}\)。该映射逐坐标保持乘法，而各投影又唯一确定它，故得到所需的唯一同态。这里没有对无限多个群元素作乘法。
\symidx{\(\prod_{i\in I}G_i\)：群的直积}

\begin{definition}\label{b1:def:restricted-product}
设 \((G_i)_{i\in I}\) 为一族群，\(g=(g_i)_{i\in I}\) 为直积中的元素，其中 \(g_i\in G_i\)。它的支集为
\(\operatorname{supp}(g)=\{\,i\in I\mid g_i\neq1\,\}\)。
所有支集有限的元素组成的集合称为\term[youxian-zhiji-zhiji]{有限支集直积}[restricted direct product with finite support]，即
\[
\Bigl\{\,(g_i)_i\in\prod_{i\in I}G_i\mid
\operatorname{supp}(g)\;\text{有限}\,\Bigr\}.
\]
其运算仍为逐坐标相乘；条件只限制哪些元素允许出现，不改变各坐标的乘法。
\end{definition}
\symidx{\(\operatorname{supp}(g)\)：直积元素的非单位坐标支集}

这个有限支集构造在英文中也常称 weak direct product。本节所用的 restricted direct product 专指相对于各因子单位子群的情形；更一般的同名构造会先指定子群 \(U_i\leq G_i\)，只要求除有限多个指标外都有 \(g_i\in U_i\)，而不要求这些坐标为单位。

\ProseParagraphBreak
它是直积的子群，因为
\(\operatorname{supp}(gh)\subseteq\operatorname{supp}(g)\cup\operatorname{supp}(h)\)，而 \(g^{-1}\) 的支集等于 \(g\) 的支集。有限指标集上它与直积相同，无限指标集上则通常不同。例如取每个 \(G_i=C_2=\{1,a_i\}\)，族 \((a_i)_i\) 属于直积；若 \(I\) 无限，它不属于有限支集直积。

\ProseParagraphBreak
有限支集直积由各坐标子群生成：任一元素只有有限个非单位坐标，故可写成这些坐标嵌入元素的有限乘积。全部直积则未必如此：在前面无限多个 \(C_2\) 的例子中，元素 \((a_i)_i\) 的支集无限，不能写成有限多个坐标嵌入元素的乘积。因此，“由坐标嵌入生成”和“由坐标投影决定”是不同性质。一般群中只定义了有限多个元素的乘积，不能以形式上的无限乘积代替这里的有限表达。

\ProseParagraphBreak
另一方面，若只有有限多个因子非平凡，则每个元素的支集都包含在这些因子的指标集合内。即使 \(I\) 无限，有限支集直积仍等于全部直积；两者的区别取决于非平凡坐标，而不只取决于指标集是否无限。

\ProseParagraphBreak
坐标子群未必生成整个无限直积，映出的同态也就未必由坐标上的限制唯一确定。用加法记号，取
\[
P=\prod_{n\geq1}\mathbb Z,\qquad
S=\Bigl\{\,(a_n)_{n\geq1}\in P\mid\{\,n\geq1\mid a_n\ne0\,\}\;\text{有限}\,\Bigr\}.
\]
这里 \(S\) 是有限支集子群。由于 \(P\) 交换，\(S\trianglelefteq P\)，所以商映射和零同态
\[
q:P\longrightarrow P/S,\qquad q(x)=x+S,
\qquad 0:P\longrightarrow P/S,\qquad 0(x)=S
\]
都有定义。每个坐标嵌入 \(\iota_n:\mathbb Z\rightarrow P\) 的像都包含在 \(S\) 中，因此 \(q\iota_n=0\iota_n=0\)。但常数序列 \(v=(1,1,\ldots)\) 的支集无限，不属于 \(S\)，故
\[
q(v)=v+S\ne S=0(v).
\]
这两个不同的同态在全部坐标子群上的限制相同，所以从无限直积映出的同态未必由这些限制决定。\cref{b1:ex:06-infinite-product-maps}进一步讨论这个商群中的元素和同态。

\subsection{交换群的直和}
\label{b1:subsec:0601:03}

\begin{definition}[交换群的直和]\label{b1:def:abelian-direct-sum}
设 \((A_i)_{i\in I}\) 是一族交换群，采用加法记号，单位写成 \(0\)。令
\[
\bigoplus_{i\in I}A_i
=\Bigl\{\,(a_i)_{i\in I}\in\prod_{i\in I}A_i\mid
\{\,i\in I\mid a_i\neq0\,\}\;\text{有限}\,\Bigr\}
\]
在这个集合上逐坐标定义加法，称所得交换群为这族群的\term[zhihe]{直和}[direct sum]。
\symidx{\(\bigoplus_{i\in I}A_i\)：交换群的直和}
零元是零族，负元逐坐标取负。坐标嵌入 \(\iota_i:A_i\rightarrow\bigoplus_{j\in I}A_j\) 将 \(a\) 放在第 \(i\) 个坐标，其余坐标均置为 \(0\)。
\end{definition}

直和正是采用加法记号的有限支集直积；各 \(A_i\) 交换，保证所得群本身交换。前面的整数序列子群 \(S\) 因而也记为 \(\bigoplus_{n\geq1}\mathbb Z\)。对于映出直和的同态，目标群 \(B\) 的交换性另有作用：它使任意两个不同坐标的像自动逐元素交换，从而不必另加像之间的交换条件。若接收群不是交换群，仍须要求这些像逐元素交换；本节最后的 \(C_2\times C_2\rightarrow S_3\) 例子将说明这一条件不能省略。

\begin{proposition}[交换群直和的泛性质]\label{b1:prop:abelian-sum-universal}
设 \((A_i)_{i\in I}\) 是一族交换群，\(B\) 是交换群，并给定群同态 \(u_i:A_i\rightarrow B\)。对每个 \(i\in I\)，令坐标嵌入
\[
\iota_i:A_i\longrightarrow\bigoplus_{j\in I}A_j
\]
把 \(a\in A_i\) 放在第 \(i\) 个坐标，并将其余坐标置为 \(0\)。则有唯一群同态
\[
u:\bigoplus_iA_i\longrightarrow B,\qquad
u\bigl((a_i)_i\bigr)=\sum_i u_i(a_i),
\]
使 \(u\iota_i=u_i\) 对全部 \(i\) 成立。式中的和总是有限和，空和按 \(0_B\) 解释。
\end{proposition}
\begin{proof}
有限支集使公式有意义。对 \(a,a'\) 取其支集的有限并，在这个共同有限集合上使用各 \(u_i\) 的可加性和 \(B\) 的交换性，可得 \(u(a+a')=u(a)+u(a')\)。由于各 \(u_j(0)=0\)，该公式还给出 \(u\iota_i=u_i\)，这证明存在性。

若 \(v\) 也满足 \(v\iota_i=u_i\)，任取有限支集元素 \(a\)，写成有限和 \(a=\sum_{i\in\operatorname{supp}(a)}\iota_i(a_i)\)。可加性便强制
\[
v(a)=\sum_{i\in\operatorname{supp}(a)}v\bigl(\iota_i(a_i)\bigr)
=\sum_{i\in\operatorname{supp}(a)}u_i(a_i)=u(a),
\]
故所构造的同态唯一。

\end{proof}

因此在交换群范围内，直和是余积：对每个接收群给出从各因子出发的同态，能唯一合并成一个同态。有限个交换群的直和与直积是同一个群，但体现两种不同方向的泛性质。无限情形中，\(\bigoplus_{n\geq1}\mathbb Z\) 的元素是有限支集整数序列，\(\prod_{n\geq1}\mathbb Z\) 则允许任意整数序列，二者不再能混用。

\ProseParagraphBreak
空指标集的情形可一并说明：只有一个空族，因此空直积和空有限支集直积都是平凡群，交换群的空直和是零群。任意群到平凡群只有一个同态，从零群到任意交换群也只有一个同态，因而上述两种方向的泛性质都包括这个情形。

\subsection{有限支集直积与群的余积}
\label{b1:subsec:0601:04}

在所有群的范围内，有限支集直积不具有上述不加限制的余积泛性质。障碍不是无限指标集，而是不同坐标之间被强制加入的交换关系。

\ProseParagraphBreak
例如有两个从 \(C_2\) 到 \(S_3\) 的同态，分别把生成元送到 \((12)\) 和 \((23)\)。若它们能延拓为 \(C_2\times C_2\rightarrow S_3\) 的同态，那么两坐标的生成元在原群中交换，其像也应交换。然而
\[
(12)(23)=(123),\qquad(23)(12)=(132),
\]
两者不同，所以延拓不存在。这已在只有两个因子的有限情形否定了群余积的泛性质。

\begin{proposition}[从有限支集直积出发的同态]\label{b1:prop:finite-support-product-maps}
设 \((G_i)_{i\in I}\) 为一族群，并令
\[
D
=\Bigl\{\,(g_i)_{i\in I}\in\prod_{i\in I}G_i
\mid \operatorname{supp}(g)\;\text{有限}\,\Bigr\}
\]
为其有限支集直积。对每个 \(i\in I\)，令 \(\iota_i:G_i\rightarrow D\) 为第 \(i\) 个坐标嵌入。给定群 \(K\) 以及群同态 \(u_i:G_i\rightarrow K\)，存在群同态
\[
u:D\longrightarrow K
\]
满足 \(u\iota_i=u_i\) 对每个 \(i\in I\) 成立，当且仅当不同 \(u_i\) 的像逐元素交换，即
\[
u_i(x)u_j(y)=u_j(y)u_i(x)
\quad(i\neq j,\ x\in G_i,\ y\in G_j).
\]
存在时，该同态唯一；对每个 \(g=(g_i)_{i\in I}\in D\)，它由
\[
u(g)=\prod_{i\in\operatorname{supp}(g)}u_i(g_i)
\]
给出，其中右端是与排列次序无关的有限乘积，空乘积按 \(1_K\) 解释。
\end{proposition}
\begin{proof}
若这样的 \(u\) 存在，则对 \(i\neq j\)、\(x\in G_i\)、\(y\in G_j\)，坐标嵌入满足
\[
\iota_i(x)\iota_j(y)=\iota_j(y)\iota_i(x).
\]
对等式两端应用 \(u\)，便得到所述交换条件。

反过来，设交换条件成立。对 \(g\in D\)，按命题中的公式定义 \(u(g)\)。它只涉及有限多个非单位坐标，不同坐标的像又逐元素交换，所以有限乘积与排列次序无关。若 \(g,h\in D\)，令 \(F=\operatorname{supp}(g)\cup\operatorname{supp}(h)\)。当 \(F=\varnothing\) 时，\(g=h=1_D\)，同态等式显然成立；否则把 \(g,h\) 看作有限直积 \(\prod_{i\in F}G_i\) 的元素。所给公式在这个有限直积上正是\cref{b1:prop:finite-product-maps}构造的同态，因此 \(u(gh)=u(g)u(h)\)。公式也直接给出 \(u\iota_i=u_i\)。

最后，每个 \(g\in D\) 都可写成有限乘积
\[
g=\prod_{i\in\operatorname{supp}(g)}\iota_i(g_i).
\]
任何在各坐标上限制为 \(u_i\) 的同态都必须把它送到命题所给的有限乘积，故 \(u\) 唯一。
\end{proof}

真正的群余积不应事先要求不同因子的像交换，它将在\secref{b1:sec:0803}中由自由积构造。后文对非交换群使用“有限支集直积”，只表示本节的具体子群，不把它当作交换群直和泛性质的无条件推广。
